\chapter{Как мы это узнали}\label{ch:kak}

Эта глава --- рецепт. По ней можно повторить разбор на своём дроне и понять, где в нём доверять, а где перепроверять.

\section{Хронология открытий}

\begin{figure}[tbp]
\centering
\begin{tikzpicture}[
  ev/.style={draw=rule, fill=paper, rounded corners=2pt, text width=36mm, minimum height=14mm, align=center, font=\sffamily\scriptsize, inner sep=3pt},
  hot/.style={ev, draw=accent, fill=accentsoft},
  ar/.style={-{Stealth[length=2mm]}, color=accentdark, line width=0.8pt}
]
\node[ev] (a) {Образы Wing: Linux, \texttt{pilot}, нет детектора};
\node[ev, right=7mm of a] (b) {Ядро C10 с символами: формула метки};
\node[ev, right=7mm of b] (c) {Приложения: протокол и команды};
\node[hot, below=10mm of a] (d) {Модуль V831 по USB: детектор найден};
\node[hot, right=7mm of d] (e) {Кадр модуля снят в памяти программы};
\node[hot, right=7mm of e] (f) {Дрон по USB: приём кадра вживую};
\draw[ar] (a)--(b); \draw[ar] (b)--(c);
\draw[ar] (c.south) -- ++(0,-5mm) -| (d.north);
\draw[ar] (d)--(e); \draw[ar] (e)--(f);
\end{tikzpicture}
\caption{Порядок открытий. Закрашенные шаги дали ответы, которых не дали файлы.}\label{fig:timeline}
\end{figure}

\section{Шаг 1. Файлы производителя}

Из образа дрона достаются корневые системы (\texttt{squashfs}), из них --- программы. Для бинарников без символов используем дизассемблер (capstone). Главные приёмы:

\begin{itemize}
  \item Строки: по ним сразу видны имена функций, форматы журнала, пути устройств.
  \item Ссылки на строки: функция, использующая строку «\texttt{use int\_dev0}», --- место, где включается этот канал.
  \item Ядро C10 с символами служит картой: находим там функцию по имени, затем её двойника в голом бинарнике.
\end{itemize}

\section{Шаг 2. Приложения}

Мобильные приложения (APK) содержат классы, построители кадров и собственные библиотеки. Для разбора использовались androguard (байткод) и дизассемблер для ARM64. Так мы получили команды Wing и раскладку телеметрии. \pr{}

\section{Шаг 3. Само устройство}

Файлы дают много, но не всё. Остальное дают сами устройства.

\begin{enumerate}
  \item \textbf{ADB.} Если устройство показывается компьютеру как \texttt{18D1:0002} и ADB его видит, есть оболочка. Команда \texttt{adb pull} копирует файлы. Подробнее в таблице ниже.
  \item \textbf{Память процесса.} Читая \texttt{/proc/PID/mem} (только чтение), можно увидеть буфер, который автопилот заполнил из UART. Так мы проверили приём кадра.
  \item \textbf{Счётчики.} Файл \texttt{/proc/tty/driver/} показывает число принятых и отправленных байтов каждого UART: сразу видно, говорит ли модуль с дроном.
\end{enumerate}

\begin{tabularx}{\textwidth}{@{}p{5.4cm} Y@{}}
\toprule
\hdr{Что нужно узнать} & \hdr{Где смотреть} \\ \midrule
Какая версия автопилота стоит & Размер и дата \texttt{/img\_app/pilot/pilot\_f}. \\
Говорит ли модуль с дроном & Счётчик \texttt{rx} у \texttt{ttyS0}. \\
Что принял автопилот & Память процесса, буфер приёма. \\
Что шлёт модуль & Память программы модуля (последний кадр). \\
Какие режимы и команды понимает модуль & Таблица в библиотеке \texttt{libmaix\_ckh831.so}. \\ \bottomrule
\end{tabularx}

\begin{opasno}[Правила работы]
Только чтение. Без пропеллеров. Перед любыми изменениями --- полные копии. Ни одной записи во флеш дрона и модуля. Питание дрона от аккумулятора (иначе он не загрузится). \pr{}
\end{opasno}

\section{Подводные камни}

\begin{itemize}
  \item Серийный номер USB \texttt{20080411} одинаков у дрона и модуля: это номер по умолчанию от Allwinner. Различайте устройства по \texttt{uname} и списку процессов.
  \item Старый \texttt{adbd} на дроне не умеет \texttt{exec-out}: файлы берутся через \texttt{pull}, а разделы флеша копируются во временный файл в оперативной памяти.
  \item На дроне нет \texttt{od}, \texttt{tr}, \texttt{head}, \texttt{strings}: разбор байтов делается на компьютере.
  \item Линейный дизассемблер может пропускать куски кода. Мы один раз «не нашли» приёмник кадра, а он был на месте. Если поиск пустой, ищите другим способом.
\end{itemize}

\section{Инструменты в репозитории}

Скрипты лежат в \texttt{tools/firmware}: распаковка образов и приложений, дизассемблеры (ARM, ARM64, RISC-V), разбор отладочной информации, чтение классов Android-приложений, восстановление формулы метки, чтение буфера модуля. Все документы разбора лежат в \texttt{docs/firmware}. \pr{}
